\documentclass[11pt]{article}

\usepackage[a4paper,margin=30mm]{geometry}
\usepackage{amsmath,amssymb,amsthm}
\usepackage[hidelinks]{hyperref}
\usepackage{microtype}

\hypersetup{
  pdftitle={Sextic characters and idempotents in prime-order Alexander quandle rings},
  pdfauthor={Xiaojian Zeng},
  pdfsubject={Sextic character criterion for idempotents in prime-order Alexander quandle rings},
  pdfkeywords={Alexander quandle, quandle ring, idempotent, sextic character, cyclotomic number, Jacobi sum}
}

\newtheorem{theorem}{Theorem}
\newtheorem{lemma}[theorem]{Lemma}
\newcommand{\Fp}{\mathbb F_p}
\newcommand{\eps}{\varepsilon}

\title{Sextic characters and idempotents in prime-order Alexander quandle rings}
\author{Xiaojian Zeng\\
\small PBC School of Finance, Tsinghua University, Beijing 100190, China\\
\small \texttt{xiaojianzeng@foxmail.com}}
\date{}

\begin{document}
\maketitle

\begin{abstract}
For a prime $p\equiv1\;(\mathrm{mod}\ 12)$, we count parameters $a\in
\mathbb F_p\setminus\{0,1\}$ for which a multiplicative character takes the
two values $\zeta_6$ and $\zeta_6^{-1}$ at $a$ and $1-a$.  The canonical
sextic character gives a subcount $S_p$ satisfying
\[
 S_p\geq \frac{p+1-14\sqrt p}{18}.
\]
This proves existence for $p\geq194$.  Exact discrete-log certificates cover
all remaining primes $p\geq37$, while an exhaustive certificate gives
$N_{13}=0$.  By the fixed-parameter character criterion of
Elhamdadi--Ta--Virgin, it follows that for every prime $p\geq37$ with
$p\equiv1\;(\mathrm{mod}\ 12)$, some prime-order Alexander quandle has a nontrivial
idempotent in its integral quandle ring; $13$ is the unique prime congruent to
$1$ modulo $12$ without such a parameter.
\end{abstract}

\section{Statement and context}

For $a\in\Fp\setminus\{0,1\}$, let $Q_{p,a}$ be the latin Alexander
quandle with operation
\[
 x*y=ax+(1-a)y.
\]
Its integral quandle ring is the free $\mathbb Z$-module on symbols $e_x$
with the bilinear multiplication $e_xe_y=e_{x*y}$.  A nontrivial idempotent
is an element $u\neq e_x$ satisfying $u^2=u$.

Fix $\zeta_6=e^{2\pi i/6}$ and define
\[
N_p=\#\left\{a\in\Fp\setminus\{0,1\}:\begin{array}{l}
\text{there is a nontrivial multiplicative character }\chi\text{ with}\\[-2pt]
\chi(-1)=1,\quad
\{\chi(a),\chi(1-a)\}=\{\zeta_6,\zeta_6^{-1}\}
\end{array}\right\}.
\]
Elhamdadi, Ta, and Virgin prove that, for each fixed admissible pair $(p,a)$,
the character condition in this definition is equivalent to the existence of
a nontrivial idempotent in $\mathbb Z[Q_{p,a}]$; it also forces
$p\equiv1\;(\mathrm{mod}\ 12)$ \cite[Lemma~6.2 and Theorem~6.3]{ETV2026}.
Their Corollary~6.4 handles primes outside this congruence class, and their
introduction records the example $(p,a)=(37,18)$.  The result below supplies
the universal existence statement for the remaining congruence class.

\begin{theorem}\label{thm:main}
If $p\geq37$ is prime and $p\equiv1\;(\mathrm{mod}\ 12)$, then $N_p>0$.
Moreover $N_{13}=0$.  Consequently, among primes congruent to $1$ modulo
$12$, the prime $13$ is the unique one for which no parameter $a$ produces a
nontrivial idempotent in $\mathbb Z[Q_{p,a}]$.
\end{theorem}

The proof uses order-six cyclotomic numbers.  These numbers go back to the
classical study of cyclotomy \cite{Dickson1935}.  Rong, Li, and Li recall the
same definition and explicit order-six evaluations for primes
$p\equiv1\;(\mathrm{mod}\ 24)$ \cite[Section~2]{RongLiLi2023}.  Those evaluations are
not needed here and, as stated there, do not cover the second residue class
$p\equiv13\;(\mathrm{mod}\ 24)$.  A direct Jacobi-sum estimate treats both classes
uniformly.

\section{A uniform Jacobi-sum bound}

Let $p\equiv1\;(\mathrm{mod}\ 12)$ and choose a primitive root $g$ modulo $p$.
Take the sextic character $\chi$ defined by $\chi(g)=\zeta_6$, and put
\[
 C_i=g^i\langle g^6\rangle\quad(0\leq i<6),\qquad
 (i,j)_6=\#\{x\in C_i:1+x\in C_j\}.
\]
Since $(p-1)/2$ is divisible by $6$, we have $-1\in C_0$ and
$\chi(-1)=1$.  Under the substitution $x=-a$, multiplication by $-1$
preserves every $C_i$.  Hence
\[
 S_p:=(1,5)_6+(5,1)_6
\]
is exactly the number of parameters witnessed by this fixed character.  In
particular, $N_p\geq S_p$.

Extend every multiplicative character, including the principal character
$\eps$, by zero at zero.  Character orthogonality gives
\[
 1_{C_i}(t)=\frac16\sum_{r=0}^{5}\zeta_6^{-ir}\chi^r(t).
\]
For $J(\alpha,\beta)=\sum_{t\in\Fp}\alpha(t)\beta(1-t)$, expansion of the
two cyclotomic numbers after $t=-x$ gives
\begin{equation}\label{eq:expansion}
 36S_p=\sum_{r,s=0}^{5}c_{r-s}J(\chi^r,\chi^s),
 \qquad
 c_d=\zeta_6^d+\zeta_6^{-d}=(2,1,-1,-2,-1,1)_d.
\end{equation}

We use the elementary identities
\[
 J(\eps,\eps)=p-2,\qquad
 J(\eps,\alpha)=J(\alpha,\eps)=-1,\qquad
 J(\alpha,\alpha^{-1})=-\alpha(-1)=-1
\]
for nonprincipal $\alpha$.  If $\alpha$, $\beta$, and $\alpha\beta$ are all
nonprincipal, then
\begin{equation}\label{eq:jacobi-modulus}
 |J(\alpha,\beta)|=\sqrt p.
\end{equation}
For completeness, choose a nontrivial additive character and form the Gauss
sum $G(\gamma)$.  Grouping pairs of nonzero variables by their sum proves
$G(\alpha)G(\beta)=J(\alpha,\beta)G(\alpha\beta)$, while grouping by their
ratio proves $|G(\gamma)|^2=p$ for nonprincipal $\gamma$; taking absolute
values yields \eqref{eq:jacobi-modulus}.

We now partition \eqref{eq:expansion}.  The $(0,0)$ term is $2(p-2)$.
On each principal axis,
\[
 \sum_{r=1}^{5}c_r=-2,
\]
so the two axes contribute $4$.  For the five nonprincipal pairs with
$r+s\equiv0\;(\mathrm{mod}\ 6)$, the coefficient sum is
\[
 c_2+c_4+c_0+c_2+c_4=-2,
\]
and their contribution is $2$.  The remaining $20$ pairs have
$r,s,r+s\not\equiv0\;(\mathrm{mod}\ 6)$; the sum of the absolute values of their
coefficients is $28$.  Therefore
\[
 36S_p=2p+2+E_p,\qquad |E_p|\leq28\sqrt p,
\]
and consequently
\begin{equation}\label{eq:bound}
 S_p\geq\frac{p+1-14\sqrt p}{18}.
\end{equation}
For every integer $p\geq194$ the numerator is positive.  Indeed,
$(p+1)^2-196p=p^2-194p+1$ is positive at $194$ and strictly increasing
thereafter.  Thus $S_p$, which is an integer, is at least $1$ throughout the
analytic range.

\section{Exact finite certificates}

The integers congruent to $1$ modulo $12$ in $[37,194)$ are
\[
37,49,61,73,85,97,109,121,133,145,157,169,181,193.
\]
Factoring the six composite entries leaves precisely the eight primes in
Table~\ref{tab:cert}.  In each row, $g$ is a primitive root and the last two
columns give the discrete logarithms modulo $6$.  Thus the two entries are
$1$ and $5$ in some order, which is an exact sextic-character certificate.

\begin{table}[ht]
\centering
\caption{Finite certificates below the analytic cutoff.}
\label{tab:cert}
\begin{tabular}{rrrrr}
$p$ & $g$ & $a$ & $\log_g(a)\bmod6$ & $\log_g(1-a)\bmod6$\\
37  & 2 & 18 & 5 & 1\\
61  & 2 & 18 & 1 & 5\\
73  & 5 & 15 & 1 & 5\\
97  & 5 & 14 & 5 & 1\\
109 & 6 & 11 & 5 & 1\\
157 & 5 & 61 & 1 & 5\\
181 & 2 & 18 & 5 & 1\\
193 & 5 & 45 & 1 & 5\\
\end{tabular}
\end{table}

For reproducibility, the accompanying integer-only program
\texttt{jacobi\_bound\_check.py} trial-divides to obtain the complete prime
list, certifies each primitive root using the prime divisors of $p-1$, builds
the full discrete-log table, and recomputes every row.  It also verifies from
the six exact integers $c_d$ that there are $20$ generic pairs and that their
absolute coefficient mass is $28$.  No floating-point arithmetic or random
sampling is used.

This certificate architecture follows the proof-carrying finite-classification
pattern in which a released verifier reconstructs the finite search space and
checks exact integer predicates; compare the explicitly documented verifier
design in \cite{202608.1658}.  The cited application concerns finite geometric
models, not quandles or character sums, and is used here only for this
methodological comparison.

It remains to prove the boundary assertion.  With primitive root $2$ modulo
$13$, write
\[
 \chi_k(2)=e^{2\pi i k/12}.
\]
The nontrivial characters satisfying $\chi_k(-1)=1$ are exactly those with
$k=2,4,6,8,10$.  The characters with $k=4,6,8$ cannot take the values
$\zeta_6^{\pm1}$.  For $k=2$ or $10$, the required logarithm classes modulo
$6$ are $1$ and $5$.  Their elements are
\[
 C_1=\{2,11\},\qquad C_5=\{6,7\}.
\]
But
\[
 1-C_1=\{12,3\},\qquad 1-C_5=\{8,7\},
\]
so neither class contains an element whose complement lies in the other.
This exhausts all allowed nontrivial characters and proves $N_{13}=0$.
Together with \eqref{eq:bound}, Table~\ref{tab:cert}, and the
fixed-parameter equivalence of \cite{ETV2026}, this proves
Theorem~\ref{thm:main}.

\section*{Computational and AI assistance disclosure}

The coefficient partition and finite discrete-log certificates were checked
by the accompanying deterministic Python~3 program, which uses only integer
arithmetic.  A generative AI system assisted with proof organization,
development of the checker, bibliography lookup, and manuscript drafting.
The mathematical argument and exact certificates are presented above so that
they can be checked independently of that assistance.

\bibliographystyle{plain}
\bibliography{references}

\end{document}
